\newcommand{\DoNotLoadEpstopdf}{}
\documentclass[11pt,reqno]{amsart}
\usepackage[T1]{fontenc}
\usepackage{lmodern,amsmath,amssymb,amsthm,mathtools,mathrsfs,microtype}
\usepackage[margin=1.1in]{geometry}
\usepackage[hidelinks]{hyperref}
\hypersetup{pdftitle={Collision records, arithmetic, and raw local inversion in a triangular Lorentz gas},pdfauthor={Qian Qi}}
\setlength{\emergencystretch}{2em}
\allowdisplaybreaks[2]
\numberwithin{equation}{section}
\newtheorem{theorem}{Theorem}[section]
\newtheorem{lemma}[theorem]{Lemma}
\newtheorem{proposition}[theorem]{Proposition}
\newtheorem{corollary}[theorem]{Corollary}
\theoremstyle{definition}
\newtheorem{definition}[theorem]{Definition}
\newtheorem{inputstatement}[theorem]{Input statement}
\theoremstyle{remark}\newtheorem{remark}[theorem]{Remark}
\newcommand{\R}{\mathbb R}\newcommand{\Z}{\mathbb Z}
\newcommand{\one}{\mathbf1}\newcommand{\dd}{\,\mathrm d}
\newcommand{\esssup}{\operatorname*{ess\,sup}}
\newcommand{\diver}{\operatorname{div}}\newcommand{\TV}{\operatorname{TV}}
\title[Collision records and raw local inversion]{Collision records, arithmetic, and raw local inversion in a triangular Lorentz gas}
\author{Qian Qi}
\date{October 13, 2026. A2-DYN, revision 98}
\subjclass[2020]{37D50, 37A50, 60F05, 28A75}
\keywords{Lorentz gas, exact collision records, critical gradient, physical area formula, optical multiplicity, paired boundary traces}
\begin{document}
\raggedbottom
\begin{abstract}
We retain the exact-label pointwise raw local limit problem for actual
returns in a triangular Lorentz gas. A physical area formula in the two
endpoint momenta gives a quadratic small-gradient cap with an explicit
weighted optical multiplicity. Combining this cap with the inherited
spectral estimate makes the first inverse-gradient layer cake integrable
to zero at every finite collision count, with its entire error retained.
We separate the preceding residual into long near-critical targets and
short physical segments. The first is paid by the new untruncated
insertion; the second is compared to the minimum of its two actual
weighted endpoint traces. The error of that comparison is a complete
physical volume term. An explicit spectral-capacity condition would
make these errors negligible and give a critical-residual height bound
of order $(qr)^{257/256}$. That condition is not verified, and the remaining paired short-trace
height is not proved to vanish for the full family. All
original sources, labels, arithmetic classes and earlier proofs remain
unchanged; the pointwise objective is not replaced by a finite-norm law.
\end{abstract}
\maketitle
\input{journal-problem-v93}

\subsection{The finite-count estimate and the two physical mechanisms}
The new estimate acts on the actual insertion $I_m/|g|$, without a
positive lower gradient cutoff. Its leading local term is
$C b^{1/128}(|J|+B_*^{-1})$. Its complete finite-count error is
bounded by an explicit function of the weighted two-dimensional
momentum-fiber multiplicity $\mathcal K_m$ and the fixed-band decay
factor $\theta_{B_*}^m$; see Theorem~\ref{thm:v98-uncut-insertion}.
The weighted multiplicity is finite at each count, but its growth has
not been bounded at the rate needed for scalar closure.

The residual is split disjointly by total physical residence, with the
overlap assigned to the short component. On the longer component a
low capped clock forces a definite decrease of the original angular
density on a full forward interval. On the short component we compare
the target density to its two physical endpoint traces using the same
Borel segment, with no trace on an artificial target-selector graph.
Both comparisons hold at finite count before any asymptotic condition.
The conditional height consequences and their exact additional
hypothesis are stated separately, in Theorems~\ref{thm:v98-critical-certificate}
and \ref{thm:v98-paired-trace}.

\input{journal-inputs-v94}
\input{journal-residual-v98}
\input{core/235_quadratic_gradient_cap_and_uncut_insertion}
\input{core/236_disjoint_residuals_and_critical_height_certificate}
\input{core/237_paired_physical_traces_and_scalar_reduction}

\appendix
\input{appendices/v98_dependency_map}
\input{journal-references-v98}
\end{document}
